\chapter{Sequences, Series and Limits}

\begin{center}\itshape
$1-\frac12+\frac13-\frac14+\cdots=\ln2$. Rearrange the same terms as $1-\frac12-\frac14+\frac13-\frac16-\frac18+\cdots$\\
and the sum becomes $\frac12\ln2$. Infinite sums are not as innocent as finite ones.
\end{center}

\section{Sequences}

\begin{reminder}[convergence of real sequences]
\begin{itemize}
\item A monotone sequence converges if and only if it is bounded.
\item \textbf{Adjacent sequences}: if $(u_n)$ increases, $(v_n)$ decreases and $v_n-u_n\to0$, both converge to the same limit.
\item \textbf{Bolzano--Weierstrass}: every bounded sequence has a convergent subsequence.
\item \textbf{Cauchy}: $(u_n)$ converges if and only if $\sup_{p,q\ge n}|u_p-u_q|\to0$.
\end{itemize}
\end{reminder}

\begin{theorem}[Ces\`aro]\label{thm:cesaro}
If $u_n\to\ell\in\R$, then $\dfrac{u_1+\dots+u_n}n\to\ell$.
\end{theorem}
\begin{proof}
Replacing $u_n$ by $u_n-\ell$ we may assume $\ell=0$. Let $\varepsilon>0$ and $N$ with $|u_n|\le\varepsilon$ for $n>N$. Then for $n>N$,
\[
\left|\frac{u_1+\dots+u_n}n\right|\le\frac{|u_1|+\dots+|u_N|}n+\frac{n-N}n\,\varepsilon,
\]
and the first term tends to $0$, so the average is at most $2\varepsilon$ in absolute value for $n$ large.
\end{proof}

A useful corollary: if $u_{n+1}-u_n\to\ell$ then $u_n/n\to\ell$ (apply Ces\`aro to the differences).

\section{Series}

\begin{reminder}[series]
\begin{itemize}
\item For $u_n\ge0$, $\sum u_n$ converges if and only if its partial sums are bounded. Hence \textbf{comparison}: if $0\le u_n\le v_n$ and $\sum v_n$ converges, so does $\sum u_n$; and if $u_n\sim v_n$ with $v_n\ge0$, the two series have the same nature.
\item \textbf{Riemann}: $\sum\frac1{n^\alpha}$ converges if and only if $\alpha>1$.
\item \textbf{Ratio test} (d'Alembert): if $u_n>0$ and $\frac{u_{n+1}}{u_n}\to\ell$, the series converges if $\ell<1$ and diverges if $\ell>1$.
\item Absolute convergence implies convergence.
\item \textbf{Alternating series}: if $(a_n)$ decreases to $0$, $\sum(-1)^na_n$ converges, and the remainder $R_n=\sum_{k>n}(-1)^ka_k$ has the sign of its first term and satisfies $|R_n|\le a_{n+1}$.
\item \textbf{Series and integrals}: if $f\ge0$ decreases on $[1,\infty[$, then $\sum f(n)$ and $\int_1^\infty f$ have the same nature.
\end{itemize}
\end{reminder}

\begin{proposition}[Euler's constant]\label{prop:gamma}
Let $H_n=\sum_{k=1}^n\frac1k$. There is a constant $\gamma\approx0.5772$ such that $H_n=\ln n+\gamma+o(1)$.
\end{proposition}
\begin{proof}
Let $u_n=H_n-\ln n$ and $v_n=H_n-\ln(n+1)$. From $\frac1{n+1}\le\int_n^{n+1}\frac{\mathrm dt}t\le\frac1n$:
\[
u_{n+1}-u_n=\frac1{n+1}-\ln\frac{n+1}n\le0,\qquad v_{n+1}-v_n=\frac1{n+1}-\ln\frac{n+2}{n+1}\ge0,
\]
and $u_n-v_n=\ln\left(1+\frac1n\right)\to0$. The sequences are adjacent and converge to a common limit $\gamma$.
\end{proof}

\begin{theorem}[Stirling]\label{thm:stirling}
$n!\sim\sqrt{2\pi n}\left(\dfrac ne\right)^n$.
\end{theorem}
Exercise~\ref{ex:stirling} proves that $n!\sim C\sqrt n\left(\frac ne\right)^n$ for some constant $C>0$; Theorem~\ref{thm:stirlingC} in the next chapter identifies $C=\sqrt{2\pi}$ with Wallis' integrals.

\section{Exercises}

\begin{exercise}[$\star$]\label{ex:nature}
Find the nature of $\sum\frac{n!}{n^n}$, and compute $\sum_{n\ge1}\frac{n^2}{2^n}$. \emph{Hint for the second: differentiate $\sum x^n$.}
\end{exercise}

\begin{exercise}[$\star$]\label{ex:bertrand}
Show that $\sum_{n\ge2}\frac1{n\ln n}$ diverges and that $\sum_{n\ge2}\frac1{n\ln^2n}$ converges.
\end{exercise}

\begin{exercise}[$\star\star$]\label{ex:ln2}
Show that $\displaystyle\sum_{n=1}^\infty\frac{(-1)^{n+1}}n=\ln2$. \emph{Hint: write the partial sum $S_{2n}$ with harmonic numbers.}
\end{exercise}

\begin{exercise}[$\star\star$]\label{ex:sin}
Let $u_0\in\left]0,\frac\pi2\right]$ and $u_{n+1}=\sin u_n$.
\begin{enumerate}[label=(\alph*)]
\item Show that $u_n\to0$.
\item Show that $\frac1{u_{n+1}^2}-\frac1{u_n^2}\to\frac13$, and deduce that $u_n\sim\sqrt{\frac3n}$.
\end{enumerate}
\end{exercise}

\begin{exercise}[$\star\star$]\label{ex:kempner}
\emph{Kempner's series.} Show that the sum of $\frac1n$ over the integers $n\ge1$ whose decimal expansion contains no digit $9$ converges, and is less than $80$.
\end{exercise}

\begin{exercise}[$\star\star\star$]\label{ex:rearr}
\emph{Rearranging the alternating harmonic series.} Let $p,q\ge1$. Rearrange $\sum\frac{(-1)^{n+1}}n$ by taking, in turn, the next $p$ positive terms
$\frac11,\frac13,\frac15,\dots$, then the next $q$ negative terms $-\frac12,-\frac14,\dots$, and so on. Show that the new series converges to
\[
\ln2+\frac12\ln\frac pq.
\]
For $p=1$, $q=2$ this is the $\frac12\ln2$ of the opening.
\end{exercise}

\begin{exercise}[$\star\star\star$]\label{ex:stirling}
\emph{Stirling, up to a constant.} Let $d_n=\ln n!-\left(n+\frac12\right)\ln n+n$.
\begin{enumerate}[label=(\alph*)]
\item Show that $d_n-d_{n+1}=\left(n+\frac12\right)\ln\left(1+\frac1n\right)-1=\frac1{12n^2}+O\left(\frac1{n^3}\right)$.
\item Deduce that $(d_n)$ converges, and that $n!\sim C\sqrt n\left(\frac ne\right)^n$ for some $C>0$.
\end{enumerate}
\end{exercise}

\begin{beauty}
Exercise~\ref{ex:rearr} is a concrete case of Riemann's theorem: a series that converges but not absolutely can be rearranged to converge to
\emph{any} real number. Absolute convergence is exactly what makes infinite sums behave like finite ones, a theme that returns in Tome~3
(expectations) and Tome~4 (Fubini). Euler's $\gamma$ is still mysterious: nobody knows whether it is irrational.
\end{beauty}

\section{Solutions}

\begin{solution}{ex:nature}
With $u_n=\frac{n!}{n^n}$, $\frac{u_{n+1}}{u_n}=\frac{(n+1)\,n^n}{(n+1)^{n+1}}=\left(1+\frac1n\right)^{-n}\to\frac1e<1$: the series converges.
For $|x|<1$, differentiating $\sum x^n=\frac1{1-x}$ and multiplying by $x$ gives $\sum nx^n=\frac x{(1-x)^2}$; doing it again,
$\sum n^2x^n=\frac{x(1+x)}{(1-x)^3}$. At $x=\frac12$: $\frac{\frac12\cdot\frac32}{\frac18}=6$.
\end{solution}

\begin{solution}{ex:bertrand}
Both $t\mapsto\frac1{t\ln t}$ and $t\mapsto\frac1{t\ln^2t}$ are positive and decreasing on $[2,\infty[$, with antiderivatives $\ln\ln t\to\infty$ and
$-\frac1{\ln t}\to0$. By comparison with integrals, the first series diverges and the second converges.
\end{solution}

\begin{solution}{ex:ln2}
Separating even and odd terms,
\[
S_{2n}=\sum_{k=1}^{2n}\frac1k-2\sum_{k=1}^n\frac1{2k}=H_{2n}-H_n=\ln(2n)+\gamma-\ln n-\gamma+o(1)\to\ln2
\]
by Proposition~\ref{prop:gamma}. Since $S_{2n+1}-S_{2n}=\frac1{2n+1}\to0$, $S_n\to\ln2$.
\end{solution}

\begin{solution}{ex:sin}
(a) Since $0<\sin x<x$ for $x>0$, $(u_n)$ is positive and decreasing, so it converges to some $\ell\in\left[0,\frac\pi2\right]$ with $\sin\ell=\ell$, hence $\ell=0$.

(b) As $x\to0$, $\sin^2x=x^2-\frac{x^4}3+O(x^6)$, so
\[
\frac1{\sin^2x}-\frac1{x^2}=\frac{x^2-\sin^2x}{x^2\sin^2x}=\frac{\frac{x^4}3+O(x^6)}{x^4+O(x^6)}\to\frac13.
\]
With $x=u_n\to0$ this gives the first claim. By the corollary of Ces\`aro, $\frac1{nu_n^2}\to\frac13$, that is $u_n\sim\sqrt{\frac3n}$.
\end{solution}

\begin{solution}{ex:kempner}
There are $8\cdot9^{k-1}$ integers with $k$ digits and no digit $9$ (the first digit has $8$ choices, the others $9$), and each is at least $10^{k-1}$.
The terms are positive, so the partial sums are bounded by
\[
\sum_{k=1}^\infty\frac{8\cdot9^{k-1}}{10^{k-1}}=\frac8{1-\frac9{10}}=80.
\]
(The true value is about $22.92$.)
\end{solution}

\begin{solution}{ex:rearr}
Write $O_m=\sum_{k=1}^m\frac1{2k-1}=H_{2m}-\frac12H_m$ and $E_m=\sum_{k=1}^m\frac1{2k}=\frac12H_m$. By Proposition~\ref{prop:gamma},
\[
O_m=\ln2+\frac12\ln m+\frac\gamma2+o(1),\qquad E_m=\frac12\ln m+\frac\gamma2+o(1).
\]
After $n$ complete blocks, the partial sum is
$O_{pn}-E_{qn}=\ln2+\frac12\ln\frac{pn}{qn}+o(1)\to\ln2+\frac12\ln\frac pq$.
A partial sum that stops inside a block differs from one taken at the end of a block by at most $p+q$ terms, each tending to $0$, so the whole sequence of
partial sums has the same limit.
\end{solution}

\begin{solution}{ex:stirling}
(a) $d_n-d_{n+1}=-\ln(n+1)-\left(n+\frac12\right)\ln n+\left(n+\frac32\right)\ln(n+1)-1=\left(n+\frac12\right)\ln\left(1+\frac1n\right)-1$. With
$\ln(1+x)=x-\frac{x^2}2+\frac{x^3}3+O(x^4)$ at $x=\frac1n$:
\[
\left(n+\frac12\right)\left(\frac1n-\frac1{2n^2}+\frac1{3n^3}+O\left(\frac1{n^4}\right)\right)-1
=\left(\frac13-\frac14\right)\frac1{n^2}+O\left(\frac1{n^3}\right)=\frac1{12n^2}+O\left(\frac1{n^3}\right).
\]
(b) So $\sum(d_n-d_{n+1})$ converges absolutely by comparison with $\sum\frac1{n^2}$, and $d_n\to d\in\R$. Exponentiating,
$\frac{n!}{\sqrt n\,(n/e)^n}=e^{d_n}\to e^d=C>0$.
\end{solution}
